1: % Reference entries, C (AN-11).  Law 1: os/iword/app/face_a.c prompt_key,
2: % format_key, block_key, retrieve_key (Cancel F1 / Esc in every prompt),
3: % paint_status (POS with Caps Lock), cursor_key (the arrows, End and the
4: % Home prefix), os/iword/core/edit.c (the unit walk: a code is a cursor
5: % position of its own; iw_ed_line_home / line_start / doc_top / doc_end /
6: % to_line), codes.c (the "Delete [code]? No (Yes)" question, the partner
7: % deleted with an attribute code; codes InitechWord does not create are
8: % stepped over).  Shots: ../wordperfect51-decomp mint MINT11, MINT27.
9: % GATES: test-pubs-iword test-iword-codes test-iword-core test-iword-face test-iword-reveal-codes
10: \Gates{test-pubs-iword test-iword-codes test-iword-core test-iword-face test-iword-reveal-codes}
11: \begin{Feature}{Cancel}
12: Cancel backs out of a menu, prompt or question without doing anything.  You
13: can press it whenever the status line shows a menu or asks you something.
14: On the editing screen, with nothing to cancel, the same key restores
15: deleted text (see \InRef{Undelete}).
16: 
17: To cancel,
18: \begin{WPSteps}
19: \item Press \FK{Cancel}{F1}
20: \Or
21: Press \K{Esc}.
22: \end{WPSteps}
23: \Documents{wpkey}{F1}
24: \Documents{wpkey}{Esc}
25: 
26: \Notes
27: \Topic{Block}
28: While Block is on, Cancel turns Block off and leaves the text as it was.
29: 
30: \Topic{Error Messages}
31: When the status line shows a message that begins \Scr{ERROR:}, any key --
32: Cancel included -- takes you back to the prompt where the error happened,
33: so that you can try again.  Press Cancel once more to return to the
34: document.
35: 
36: \Topic{Exit and Save}
37: Cancel at any of the questions of Exit (\Scr{Save Document?},
38: \Scr{Exit WP?}) or at the name prompt of Save or Retrieve returns you to the
39: document as it was.
40: 
41: \Topic{Format Screens}
42: On the Format screens Cancel goes back one screen at a time: from a
43: measurement you are typing to the Format: Line screen (leaving the setting
44: as it was), from Format: Line to Format, and from Format to the document.
45: \FK{Exit}{F7} on the Format: Line screen returns straight to the document.
46: 
47: \Topic{Moving Text}
48: At \Scr{Move cursor; press Enter to retrieve.} Cancel leaves the text where
49: it is now: text you moved stays out of the document, but you can still put
50: it back with Move's Retrieve (see \InRef{Move, Retrieve}).  When InitechWord
51: has highlighted a sentence, paragraph or page for Move, Cancel removes the
52: highlight.
53: 
54: \Topic{On the Editing Screen}
55: With no menu or prompt on the screen, \FK{Cancel}{F1} starts Undelete: it
56: shows the last text you deleted and offers to restore it (see
57: \InRef{Undelete}).  \K{Esc} does nothing on the editing screen.
58: 
59: \Topic{Search and Replace}
60: Cancel at the \Scr{Srch:} prompt leaves the cursor where it was.  At the
61: questions and strings of Replace it leaves Replace; at
62: \Scr{Replace? Yes (No)} it stops replacing.  See \InRef{Search} and
63: \InRef{Replace}.
64: 
65: \Topic{Setup}
66: Cancel (or Esc) on any Setup screen returns to the document; Exit (F7)
67: goes back one screen.  See \InRef{Setup}.
68: 
69: \SeeAlso{Block; Exit; Format; Move, Retrieve; Undelete}
70: \end{Feature}
71: 
72: \begin{Feature}{Caps Lock}
73: The Caps Lock key of your keyboard works in InitechWord as in every FLAIR
74: application.  While Caps Lock is on, the status line shows the word
75: \Scr{Pos} in capitals, \Scr{POS}, as a reminder.
76: 
77: To see whether Caps Lock is on,
78: \begin{WPSteps}
79: \item Look at the status line: \Scr{POS} means Caps Lock is on, \Scr{Pos}
80: that it is off.
81: \end{WPSteps}
82: 
83: \BeginReplay{ref-caps}
84: \CapsLock{on}\Keys{End}%
85: With Caps Lock on, a new document shows this status line:
86: \WordScreen[rows=1]{ref-caps}{The status line with Caps Lock on}
87: \CapsLock{off}%
88: \EndReplay
89: 
90: \SeeAlso{Status Line}
91: \end{Feature}
92: 
93: \begin{Feature}{Codes}
94: InitechWord keeps everything that formats a document as hidden codes in the
95: text.  When you press Bold, a hard return, a tab or a margin setting, a code
96: goes into the document at the cursor.  The editing screen shows what the
97: codes do, but not the codes; Reveal Codes shows them (see \InRef{Reveal
98: Codes}).  Appendix B lists every code and its meaning.
99: 
100: Most features put one code in the document.  Appearances and sizes put a
101: pair of codes, one where the attribute begins and one where it ends: the
102: \emph{on} code is named in capitals (\Code{[BOLD]}), the \emph{off} code in
103: lower case (\Code{[bold]}).
104: 
105: To delete a code,
106: \begin{WPSteps}
107: \item Press \FK{Reveal Codes}{Alt-F3} so that you can see the code.
108: \item Move the cursor onto the code (the code is highlighted) and press
109: \K{Delete}
110: \Or
111: Move the cursor just after the code and press \K{Backspace}.
112: \item Press \FK{Reveal Codes}{Alt-F3} again to return to the editing screen.
113: \end{WPSteps}
114: \begin{WPNote}Deleting either code of a pair deletes both.\end{WPNote}
115: 
116: \Notes
117: \Topic{Codes InitechWord Keeps but Does Not Make}
118: A document written by WordPerfect 5.1 can hold codes that InitechWord does
119: not put in documents itself: \Code{[Center]}, \Code{[Flsh Rgt]} and
120: \Code{[ ]} (a hard space), and codes Reveal Codes shows as \Code{[?]}.
121: InitechWord keeps them where they are and saves them unchanged.  Backspace
122: and Delete step over them instead of deleting them, in Reveal Codes as on
123: the editing screen.  The editing screen draws the text after a
124: \Code{[Center]} or \Code{[Flsh Rgt]} at the left margin.
125: 
126: \Topic{Cursor on a Code}
127: Each code is a place the cursor can stop.  On the editing screen a code takes
128: no room, so pressing \K{Right} across one may seem to leave the cursor where
129: it was; in Reveal Codes you see it move onto the code.
130: 
131: \Topic{Deleting Codes on the Editing Screen}
132: If Backspace or Delete reaches a code of an appearance or size, a margin
133: setting, a justification or a line spacing while Reveal Codes is off,
134: InitechWord does not delete it at once.  It asks, naming the code as Reveal
135: Codes shows it:
136: \StatusLine{Delete [BOLD]? No (Yes)}
137: Type \Typed{y} to delete the code (and, for a pair, its partner).  Any other
138: key answers No: nothing is deleted, and the cursor moves past the code -- to
139: before it for Backspace, after it for Delete -- so that the next Backspace
140: or Delete reaches the text.  Tabs, indents, hard returns and hard pages are
141: deleted without a question.
142: 
143: \BeginReplay{ref-codes}
144: For example, type \TypeRun{Read the }(with the space), press \Press{Bold}{F6}
145: and type \TypeRun{new}.  The cursor is just before \Code{[bold]}.  Press
146: \PressKey{Delete}:
147: \WordScreen[rows=1]{ref-codes-1}{Delete meets the off code of bold}
148: Type \TypeName{y}: both bold codes are deleted and \textit{new} is ordinary
149: text again.
150: \WordScreen[rows=1]{ref-codes-2}{The bold codes deleted}
151: \EndReplay
152: 
153: \Topic{Order of Codes}
154: When you choose a format at the cursor, InitechWord puts its code there.
155: Margin settings and justification take effect at the start of a line, so if
156: the cursor is not at the left margin InitechWord first ends the line there
157: (see \InRef{Margins, Left and Right} and \InRef{Justification}).
158: 
159: \SeeAlso{Appearance; Delete; Reveal Codes; Size}
160: \end{Feature}
161: 
162: \begin{Feature}{Cursor Movement}
163: The cursor keys move the cursor through the text without changing it.  The
164: cursor marks where the next character you type goes; it is the cell drawn
165: in reverse video on the editing screen.
166: 
167: To move the cursor,
168: \begin{WPSteps}
169: \item Press an arrow key to move one character left or right, or one line
170: up or down.
171: \item Press \K{End} to move to the end of the line.
172: \item Press \K{Home} before an arrow key to move farther (see the table
173: below).
174: \end{WPSteps}
175: \Documents{wpkey}{Up}
176: \Documents{wpkey}{Down}
177: \Documents{wpkey}{Left}
178: \Documents{wpkey}{Right}
179: \Documents{wpkey}{End}
180: \Documents{wpkey}{Home}
181: \Documents{wpkey}{Home,Left}
182: \Documents{wpkey}{Home,Right}
183: \Documents{wpkey}{Home,Up}
184: \Documents{wpkey}{Home,Down}
185: \Documents{wpkey}{Home,Home,Up}
186: \Documents{wpkey}{Home,Home,Down}
187: \Documents{wpkey}{Home,Home,Home,Left}
188: \Documents{wpkey}{Home,Home,Home,Up}
189: 
190: \K{Home} is a prefix key: pressing it does nothing by itself, but changes
191: what the next arrow key does.  Press Home, release it, then press the arrow.
192: \par\addvspace{4pt}
193: \noindent\begin{tabular}{@{}lp{50mm}@{}}
194: \textbf{Keys} & \textbf{Moves the cursor to}\\
195: \K{Left}, \K{Right} & the next character to the left or right\\
196: \K{Up}, \K{Down} & the line above or below\\
197: \K{End} & the end of the line\\
198: \K{Home, Left} & the start of the line (after any codes there)\\
199: \K{Home, Home, Left} & the start of the line (after any codes there)\\
200: \K{Home, Home, Home, Left} & the start of the line, before any codes there\\
201: \K{Home, Right} & the end of the line\\
202: \K{Home, Up} & the top line on the screen\\
203: \K{Home, Down} & the bottom line on the screen\\
204: \K{Home, Home, Up} & the top of the document, after the formatting codes there\\
205: \K{Home, Home, Home, Up} & the top of the document, before any codes\\
206: \K{Home, Home, Down} & the end of the document\\
207: \end{tabular}
208: \par\addvspace{4pt}
209: 
210: \Notes
211: \Topic{Codes}
212: The cursor stops on each code as it does on each character.  On the editing
213: screen a code takes no room, so a press of Left or Right that moves onto or
214: off a code seems to do nothing; with Reveal Codes on you can see where the
215: cursor is (see \InRef{Codes}).
216: 
217: \Topic{Home Twice}
218: Two Homes before Left do what one does, and so do two Homes before Right:
219: there is no farther left or right than the start and end of a line.  With Up and Down, two Homes go to the top or end of
220: the whole document.
221: 
222: \Topic{Home, Up and Home, Down at the Edge of the Screen}
223: If the cursor is already on the top line of the screen, \K{Home, Up} moves
224: it up one screen; on the bottom line, \K{Home, Down} moves it down one
225: screen.
226: 
227: \Topic{Moving Up and Down}
228: Up and Down keep the cursor in the same column as it moves from line to
229: line.  On a shorter line the cursor goes to the end of the line, and returns
230: to its column on the next line long enough to have one.
231: 
232: \Topic{Shift, Ctrl and Alt}
233: The cursor keys do the same with Shift, Ctrl or Alt held as without
234: (Appendix F).
235: 
236: \Topic{The Screen}
237: The editing screen shows 23 lines.  When the cursor moves past the top or
238: bottom of the screen, the screen moves with it.
239: 
240: \SeeAlso{Block; Codes; Status Line}
241: \end{Feature}
